\exercisesection{}

\begin{enumerate}
\def\labelenumi{\arabic{enumi})}
\item
  Let G be a locally compact group. Let \(\pi\) be a unitary representation of G in a complex Hilbert space E and \(x \in \mathrm{E}\) . We have \(\pi(g)x = x\) for all \(g \in \mathrm{G}\) if and only if the positive-definite function \(g \mapsto \langle x | \pi(g)x \rangle\) is constant.
\item
  Let G be a locally compact group. For every positive-definite function \(\varphi\) on G, we denote by \((\pi_{\varphi}, x_{\varphi})\) a cyclic Hilbert representation of \(\varphi\) .
\end{enumerate}

\begin{enumerate}
\def\labelenumi{\alph{enumi})}
\item
  Let \(\varphi_1, \varphi_2\) and \(\varphi_3\) be positive-definite functions on G. If \(\varphi_3 = \varphi_1 + \varphi_2\) , then there exists an injective G-morphism from \(\pi_{\varphi_1}\) into \(\pi_{\varphi_3}\) .
\item
  If \(\varphi\) is a constant nonzero positive-definite function, then \(\pi_{\varphi}\) is trivial of dimension 1.
\item
  For every unitary representation \(\pi\) of G in a Hilbert space E, and for every \(x \in \mathrm{E}\) , there exists an injective G-morphism from \(\pi_{\varphi}\) into \(\pi\) , where \(\varphi(g) = \langle x \mid \pi(g)x \rangle\) .
\item
  Let H be a closed subgroup of G and let \(\varphi\) be a positive-definite function on G. Denote by \(\psi\) the restriction of \(\varphi\) to H; this is a positive-definite function on H, and there exists an injective H-morphism from \(\pi_{\psi}\) into \(\operatorname{Res}_{\mathrm{H}}^{\mathrm{H}}(\pi)\) .
\end{enumerate}

\begin{enumerate}
\def\labelenumi{\arabic{enumi})}
\setcounter{enumi}{2}
\tightlist
\item
  \begin{enumerate}
  \def\labelenumii{\alph{enumii})}
  \tightlist
  \item
    Give an example of an involutive Banach algebra admitting a positive linear form which is not continuous.
  \end{enumerate}
\end{enumerate}

\begin{enumerate}
\def\labelenumi{\alph{enumi})}
\setcounter{enumi}{1}
\tightlist
\item
  Give an example of an involutive Banach algebra A, a continuous positive linear form \(\lambda\) on A, and Hilbert realizations \((\mathrm{E}_{1}, e_{1}, \varphi_{1})\) and \((\mathrm{E}_{2}, e_{2}, \varphi_{2})\) of \(\lambda\) such that \((\mathrm{E}_{1}, e_{1}, \varphi_{1})\) is a cyclic Hilbert realization and such that the unique linear map u of prop. 6 of V, p.~440 does not satisfy \(u(e_{1}) = e_{2}\) .
\end{enumerate}

\begin{enumerate}
\def\labelenumi{\arabic{enumi})}
\setcounter{enumi}{3}
\tightlist
\item
  Let A and B be star algebras. A linear map u from A to B is said to be positive if \(u(\mathrm{A}_{+}) \subset \mathrm{B}_{+}\) , and completely positive if, for every integer \(k \in N\) , the map \(u_{k} = 1_{C^{k}} \otimes u: C^{k} \otimes A \to C^{k} \otimes B\) is positive.
\end{enumerate}

\begin{enumerate}
\def\labelenumi{\alph{enumi})}
\tightlist
\item
  Let A be a unital star algebra, E and F complex Hilbert spaces, v a linear map from E to F, and \(\varphi\) a unital morphism of involutive algebras from A to \(\mathcal{L}(\mathrm{F})\) such that \(\|\varphi\|^{2} = \|\varphi(1)\|\) . The map \(u: A \to \mathcal{L}(E)\) defined by \(u(a) = v^{*} \circ \varphi(a) \circ v\) is completely positive.
\end{enumerate}

The goal of the exercise is to prove the converse of this assertion (“Stinespring theorem”). Let A be a unital star algebra, E a complex Hilbert space, and \(u \colon \mathrm{A} \to \mathcal{L}(\mathrm{E})\) a completely positive linear map.

\begin{enumerate}
\def\labelenumi{\alph{enumi})}
\setcounter{enumi}{1}
\tightlist
\item
  Let \(F = A \otimes E\) . There exists a positive sesquilinear form q on F such that
\end{enumerate}

\[
q (a \otimes x, b \otimes y) = \langle x | \varphi (a ^ {*} b) y \rangle
\]

for all \((a,b)\in \mathrm{A}^2\) and \((x,y)\in \mathrm{E}^2\) .

Let \(\widehat{\mathbf{F}}\) be the separated completion Hilbert space of \(\mathbf{F}\) for the sesquilinear form \(q\) .

\begin{enumerate}
\def\labelenumi{\alph{enumi})}
\setcounter{enumi}{2}
\tightlist
\item
  Let \(b \in A\) . The endomorphism of F such that \(a \otimes x \mapsto ba \otimes x\) for all \(a \in A\) and \(x \in E\) induces by passage to quotients an endomorphism \(\varphi(b)\) of \(\widehat{F}\) such that \(\|\varphi(b)\|\leqslant\|b\|\) ; the map \(\varphi\colon A\to\mathcal{L}(F)\) is a unital morphism of involutive algebras.
\end{enumerate}

\begin{enumerate}
\def\labelenumi{\alph{enumi})}
\setcounter{enumi}{3}
\tightlist
\item
  Let \(v\) be the linear map from E to \(\widehat{\mathrm{F}}\) such that \(v(x)\) is the class of \(1 \otimes x\) for all \(x \in \mathrm{E}\) . We have \(\|v\|^2 = \| \varphi(1)\|\) , and \(u(a) = v^* \circ \varphi(a) \circ v\) for all \(a \in \mathrm{A}\) .
\end{enumerate}

\begin{enumerate}
\def\labelenumi{\arabic{enumi})}
\setcounter{enumi}{4}
\tightlist
\item
  Let X be a separated topological space. A continuous function \(f: X \times X \to R\) is said to be of conditionally negative type if it satisfies the conditions
\end{enumerate}

\begin{enumerate}
\def\labelenumi{(\roman{enumi})}
\item
  For every \(x \in X\) , we have \(f(x, x) = 0\) ;
\item
  For all \(x\) and \(y\) in \(X\) , we have \(f(x,y) = f(y,x)\) ;
\item
  For every integer \(n \in N\) , every family \((x_{i})_{0 \leqslant i \leqslant n}\) of elements of X and every family \((t_{i})_{0 \leqslant i \leqslant n}\) of real numbers such that
\end{enumerate}

\[
\sum_ {i = 1} ^ {n} t _ {i} = 0,
\]

we have

\[
\sum_ {i = 0} ^ {n} \sum_ {j = 0} ^ {n} t _ {i} t _ {j} f (x _ {i}, x _ {j}) \leqslant 0.
\]

We denote by Noy \(_{-}\) (X) the set of functions of conditionally negative type on X × X.

\begin{enumerate}
\def\labelenumi{\alph{enumi})}
\item
  Let E be a real Hilbert space and \(\varphi\colon\mathrm{X}\to\mathrm{E}\) a continuous map. The function \(f(x,y)=\|\varphi(x)-\varphi(y)\|^{2}\) is of conditionally negative type.
\item
  Let \(f\in\mathrm{Noy}_{-}(\mathrm{X})\) and \(x_{0}\in\mathrm{X}\) . There exists a real Hilbert space E and a continuous map \(\varphi\colon\mathrm{X}\to\mathrm{E}\) such that
  \begin{enumerate}
  \def\labelenumii{\alph{enumii})}
  \tightlist
  \item
    For every \((x,y)\in \mathrm{X}\times \mathrm{X}\) , we have \(f(x,y) = \| \varphi (x) - \varphi (y)\| ^2\) ;
  \item
    The family \((\varphi(x)-\varphi(x_{0}))_{x\in\mathrm{X}}\) is total in E.
  \end{enumerate}
\end{enumerate}

We say that \((\mathrm{E},\varphi)\) is a cyclic Hilbert representation of f.

\begin{enumerate}
\def\labelenumi{\alph{enumi})}
\setcounter{enumi}{2}
\item
  Let \(f \in \mathrm{Noy}_{-}(\mathrm{X})\) and \(x_0 \in \mathrm{X}\) . Let \((\mathrm{E}, \varphi)\) and \((\mathrm{E}', \varphi')\) be cyclic Hilbert representations of \(f\) . There exists a unique affine isometry \(u: \mathrm{E} \to \mathrm{E}'\) such that \(\varphi' = u \circ \varphi\) .
\item
  The set \(\mathrm{Noy}_{-}(\mathrm{X})\) is a convex cone in \(\mathcal{C}(\mathrm{X}\times \mathrm{X};\mathbf{R})\) . For every universally positive kernel \(g\in \mathrm{Noy}_{+}(\mathrm{X})\) such that \(g\) is real-valued, the map \(f(x,y) = g(x,x) - g(x,y)\) belongs to \(\mathrm{Noy}_{-}(\mathrm{X})\) .
\item
  Let \(\mathfrak{F}\) be a filter on \(\mathrm{Noy}_{-}(\mathrm{X})\) which converges simply to \(f\colon \mathrm{X}\times \mathrm{X}\to \mathbf{R}\) . If \(f\) is continuous, then \(f\in \mathrm{Noy}_{-}(\mathrm{X})\) .
\item
  Let \(f: X \times X \to R\) be a continuous function satisfying conditions (i) and (ii) above. Let \(x_{0} \in X\). Let \(g: X \times X \to R\) be such that \(g(x, y) = f(x, x_{0}) + f(y, x_{0}) - f(x, y)\) . Then \(f \in \text{Noy}_{-}(X)\) if and only if \(g \in \text{Noy}_{+}(X)\) .
\item
  Suppose X is nonempty. Let \(f: \mathrm{X} \times \mathrm{X} \to \mathbf{R}\) be a continuous function satisfying conditions (i) and (ii) above. Then \(f \in \mathrm{Noy}_{-}(\mathrm{X})\) if and only if we have \(\exp(-tf) \in \mathrm{Noy}_{+}(\mathrm{X})\) for all \(t \in \mathbf{R}_{+}\) (“Schoenberg theorem”). (To prove that \(\exp(-f)\) is of positive type when \(f \in \mathrm{Noy}_{-}(\mathrm{X})\) , fix \(x_0 \in \mathrm{X}\) ; show that
\end{enumerate}

\[
g (x, y) = \exp (f (x, x _ {0}) + f (y, x _ {0}) - f (x, y))
\]

and

\[
h (x, y) = \exp (- f (x, x _ {0})) \exp (- f (y, x _ {0}))
\]

belong to Noy \(_{+}\) (X).)

If G is a topological group, we say that a function \(\varphi \in \mathcal{C}(\mathrm{G})\) is of negative type if the function \(f\) on \(\mathrm{G} \times \mathrm{G}\) defined by \(f(g,h) = \varphi(h^{-1}g)\) is a conditionally negative-definite kernel.

\begin{enumerate}
\def\labelenumi{\arabic{enumi})}
\setcounter{enumi}{5}
\tightlist
\item
  Let G be a discrete topological group. Let A be the involutive subalgebra of \(\mathcal{L}(\ell^{2}(\mathrm{G}))\) generated by the image of the left regular representation \(\lambda\) . This is a unital star algebra.
\end{enumerate}

For every function of positive type \(\varphi\in\mathrm{Pos}(\mathrm{G})\) , there exists a unique completely positive linear map \(u_{\varphi}\) from A to A (cf.~exercise 4) such that \(u_{\varphi}(\boldsymbol{\lambda}(g))=\varphi(g)\boldsymbol{\lambda}(g)\) for all \(g\in\mathrm{G}\) (Use a cyclic Hilbertian representation of \(\varphi\) .)

\begin{enumerate}
\def\labelenumi{\arabic{enumi})}
\setcounter{enumi}{6}
\tightlist
\item
  Let I be a finite set and \(G = F(I)\) the free group constructed on I (A, I, p.~84). Let A be the involutive subalgebra of \(\mathcal{L}(\ell^{2}(G))\) generated by the image of the left regular representation \(\lambda\) .
\end{enumerate}

\begin{enumerate}
\def\labelenumi{\alph{enumi})}
\item
  For every \(g \in G\) , we denote by \(|g|\) the length of the unique reduced decomposition of g (A, I, p.~146, exercise 19). The map \(d(g, h) = |h^{-1}g|\) is a metric on G.
\item
  Let \(\lambda \in \mathbf{R}_{+}^{*}\) . The function on G defined by \(\psi_{\lambda}(g) = \exp(-\lambda|g|)\) is a positive-definite function. (Show that the function \(g \mapsto |g|\) is a function of negative type, cf.~exercise 5.)
\item
  For every integer \(k \in N\) , let \(G_{k}\) be the set of \(g \in G\) such that \(|g| = k\) and let \(\varphi_{k}\) be the characteristic function of \(G_{k}\) . Let f and g be functions on G with support in \(G_{k}\) and \(G_{l}\) respectively. For every \(m \in N\) , we have \((f * g)\varphi_{m} = 0\) unless \(k + l - m\) is even and \(|k - l| \leqslant m \leqslant k + l\) ; in this case, we have
\end{enumerate}

\[
\left\| (f * g) \varphi_ {m} \right\| \leqslant \| f \| \| g \|.
\]

\begin{enumerate}
\def\labelenumi{\alph{enumi})}
\setcounter{enumi}{3}
\tightlist
\item
  Let f be a function on G with finite support. We have
\end{enumerate}

\[
\| \boldsymbol {\lambda} (f) \| \leqslant \sum_ {n \in \mathbf {N}} (n + 1) \| f \varphi_ {n} \|,
\]

and

\[
\left\| \boldsymbol {\lambda} (f) \right\| \leqslant 2 \left(\sum_ {g \in G} (1 + | g |) ^ {4} | f (g) | ^ {2}\right) ^ {1 / 2}.
\]

\begin{enumerate}
\def\labelenumi{\alph{enumi})}
\setcounter{enumi}{4}
\item
  Let \(\psi\) be a function on G such that \(M = \sup_{g \in G} |\psi(g)| (1 + |g|)^2\) is finite. There exists a unique linear map \(u_\psi\) from A to A such that \(u_\psi(\boldsymbol{\lambda}(g)) = \psi(g)\boldsymbol{\lambda}(g)\) for all \(g \in G\) . We have \(\|u_\psi\| \leqslant 2M\) .
\item
  Let X be the set of functions \(\psi\) on G with finite support such that \(\|u_{\psi}\|\leqslant1\) . There exists a filter \(\mathfrak{F}\) on X such that
\end{enumerate}

\[
\lim _ {\psi , \mathfrak {F}} u _ {\psi} = 1 _ {\mathrm{A}}
\]

in the space \(\mathcal{L}(\mathrm{A})\) endowed with the topology of simple convergence. In particular, the Banach space underlying A has the approximation property. (Consider the functions \(\psi_{\lambda,n}(g) = (\varphi_0 + \dots + \varphi_n)\psi_\lambda\) for \(\lambda \in \mathbf{R}_+^*\) and \(n \in \mathbf{N}\) , and use the preceding exercise.)

\begin{enumerate}
\def\labelenumi{\arabic{enumi})}
\setcounter{enumi}{7}
\tightlist
\item
  One identifies the group \(\mathbf{R}\) and its dual group by the map \((x,y)\mapsto\) \(\exp (i\pi xy)\) . Let \(\alpha \in \mathbf{R}_{+}\) .
\end{enumerate}

\begin{enumerate}
\def\labelenumi{\alph{enumi})}
\item
  There exists a positive measure on \(\mathbf{R}\) whose Fourier transform is the function \(f_{\alpha}(x) = \exp(-|x|^{\alpha})\) if and only if \(\alpha \in ]0,2]\) . When this is the case, one denotes by \(\mu_{\alpha}\) this measure, which is unique; it is called the stable measure of parameter \(\alpha\) .
\item
  The measure \(\mu_{\alpha}\) has total mass 1; the identity function f on R is integrable with respect to \(\mu_{\alpha}\) if and only if \(\alpha = 2\) .
\end{enumerate}

\begin{enumerate}
\def\labelenumi{\arabic{enumi})}
\setcounter{enumi}{8}
\tightlist
\item
  The goal of this exercise is to present another proof of Theorem 5 of V, p.~455. Let G be a locally compact commutative topological group, equipped with a Haar measure \(\mu\) . We denote by \(\widehat{f}\) the Fourier transform of \(f \in \mathrm{L}^1(\mathrm{G})\) , and we denote by \(\mathcal{F}(\mathrm{G}) \subset \mathcal{C}_0(\widehat{\mathrm{G}})\) the image of the Fourier transform on \(\mathrm{L}^1(\mathrm{G})\) .
\end{enumerate}

\begin{enumerate}
\def\labelenumi{\alph{enumi})}
\item
  The space \(\mathcal{F}(\mathrm{G})\) is a dense subspace of \(\mathcal{C}_0(\mathrm{G})\) .
\item
  For every \(f \in \mathrm{L}^{1}(\mathrm{G})\) , the spectral radius of f in the Banach algebra \(\mathrm{L}^{1}(\mathrm{G})\) is equal to \(\|\widehat{f}\|_{\infty}\) .
\item
  Let \(\varphi\in\mathrm{Pos}(\mathrm{G})\) . It is a uniformly continuous and bounded function on \(\mathrm{G}\) ; for every \(f\in\mathrm{L}^{1}(\mathrm{G})\) , we have
\end{enumerate}

\[
\int_ {G} \int_ {G} f (x) \overline {{f (y)}} \varphi (x - y) d (\mu \otimes \mu) (x, y) \geqslant 0.
\]

\begin{enumerate}
\def\labelenumi{\alph{enumi})}
\setcounter{enumi}{3}
\tightlist
\item
  Let \(\lambda\) be the continuous linear form on \(\mathrm{L}^{1}(\mathrm{G})\) determined by \(\varphi\) and we set \(b(f,g)=\lambda(f^{*}*g)\) for all \((f,g)\in\mathrm{L}^{1}(\mathrm{G})\times\mathrm{L}^{1}(\mathrm{G})\) . The map \(b\) is a positive hermitian form on \(\mathrm{L}^{1}(\mathrm{G})\) .
\end{enumerate}

\begin{enumerate}
\def\labelenumi{\alph{enumi})}
\setcounter{enumi}{4}
\tightlist
\item
  For every \(f \in \mathrm{L}^1(\mathrm{G})\) , we have \(|\lambda(f)|^2 \leqslant b(f, f)\) .
\end{enumerate}

\begin{enumerate}
\def\labelenumi{\alph{enumi})}
\setcounter{enumi}{5}
\tightlist
\item
  Let \(f \in \mathrm{L}^{1}(\mathrm{G})\) and \(g = f^{*} * f \in \mathcal{C}_{b}(\mathrm{G})\) . One defines \(g_{1} = g\) and \(g_{n+1} = g_{n} * g\) for every integer \(n \geqslant 1\) . For every integer \(n \geqslant 1\) , we have \(|\lambda(f)| \leqslant \|g_{n}\|_{1}^{2^{-n}}\) .
\end{enumerate}

\begin{enumerate}
\def\labelenumi{\alph{enumi})}
\setcounter{enumi}{6}
\tightlist
\item
  For every \(f \in \mathrm{L}^{1}(\mathrm{G})\) , we have \(|\lambda(f)| \leqslant \| \widehat{f} \|_{\infty}\) . There exists a continuous linear map \(\widehat{\lambda}\) on \(\mathcal{F}(\mathrm{G})\) such that \(\lambda(f) = \widehat{\lambda}(\widehat{f})\) for all \(f \in \mathrm{L}^{1}(\mathrm{G})\) .
\end{enumerate}

\begin{enumerate}
\def\labelenumi{\alph{enumi})}
\setcounter{enumi}{7}
\tightlist
\item
  There exists a bounded measure \(\nu\) on \(\widehat{G}\) such that \(\lambda(f) = \nu(\widehat{f})\) for every \(f \in \mathrm{L}^{1}(\mathrm{G})\) ; the measure \(\nu\) is positive and one has
\end{enumerate}

\[
\varphi (x) = \int_ {\widehat {\mathrm{G}}} \langle \chi , x \rangle d \nu (\chi)
\]

for all \(x \in G\) .

\begin{enumerate}
\def\labelenumi{\arabic{enumi})}
\setcounter{enumi}{9}
\tightlist
\item
  \begin{enumerate}
  \def\labelenumii{\alph{enumii})}
  \tightlist
  \item
    Let \(\mathrm{G} = \mathbf{R}\) . The set \(\operatorname{Pos}_1(\mathrm{G})\) is not compact for the weak topology.
  \end{enumerate}
\end{enumerate}

\begin{enumerate}
\def\labelenumi{\alph{enumi})}
\setcounter{enumi}{1}
\tightlist
\item
  Let G = R/Z. The weak topology on \(\operatorname{Pos}_{\leqslant1}(\mathbf{G})\) does not coincide with the topology of compact convergence.
\end{enumerate}

\begin{enumerate}
\def\labelenumi{\arabic{enumi})}
\setcounter{enumi}{10}
\tightlist
\item
  Let G be a locally compact group.
\end{enumerate}

Let \(\varrho\) and \(\pi\) be unitary representations of G. One says that \(\varrho\) is weakly contained in \(\pi\) , and one writes \(\varrho \leqslant \pi\) , if the set of diagonal matrix coefficients of \(\varrho\) is contained in the closure of the set of finite sums of diagonal matrix coefficients of \(\pi\) in the space \(\mathcal{C}(\mathrm{G})\) endowed with the topology of compact convergence.

\begin{enumerate}
\def\labelenumi{\alph{enumi})}
\item
  The relation « \(\varrho\) and \(\pi\) are unitary representations of G and \(\varrho \leqslant \pi\) » is a preorder relation. If \(\varrho'\) (resp. \(\pi'\) ) is isomorphic to \(\varrho\) (resp. \(\pi\) ) then \(\varrho \leqslant \pi\) if and only if \(\varrho' \leqslant \pi'\) .
\item
  If \(\varrho\) is isomorphic to a subrepresentation of \(\pi\) , then \(\varrho \leqslant \pi\) . If I and J are nonempty finite sets, and if \(\varrho \leqslant \pi\) , then \(\bigoplus_{i \in I} \varrho \leqslant \bigoplus_{j \in J} \pi\) .
\item
  Let X be a subset of the space E of \(\varrho\) which generates \(\pi\) . We have \(\varrho \leqslant \pi\) if and only if every matrix coefficient \(g \mapsto \langle x | \varrho(g)x \rangle\) with \(x \in X\) belongs to the closure of the set of finite sums of diagonal matrix coefficients of \(\pi\) in the space \(\mathcal{C}(\mathrm{G})\) endowed with the topology of compact convergence. (Show that the set Y of elements of E satisfying this property is a closed vector subspace of E.)
\item
  Let H be a closed subgroup of G. Let \(\pi\) and \(\varrho\) be unitary representations of G such that \(\varrho \leqslant \pi\) ; we then have \(\operatorname{Res}_{\mathrm{H}}^{\mathrm{G}}(\varrho) \leqslant \operatorname{Res}_{\mathrm{H}}^{\mathrm{G}}(\pi)\) .
\item
  Let H be a closed subgroup of G and let \(\nu\) be a Haar measure on H. Let \(\pi\) and \(\varrho\) be unitary representations of H such that \(\varrho \leqslant \pi\) ; we then have \(\operatorname{Ind}_{\mathrm{H}}^{\mathrm{G}}(\varrho) \leqslant \operatorname{Ind}_{\mathrm{H}}^{\mathrm{G}}(\pi)\) . (Let \(\varrho' = \operatorname{Ind}_{\mathrm{H}}^{\mathrm{G}}(\varrho)\) and \(\pi' = \operatorname{Ind}_{\mathrm{H}}^{\mathrm{G}}(\pi)\) ; using part (c), reduce to approximating matrix coefficients \(g \mapsto \langle f | \varrho'(g)f \rangle\) for f of the form
\end{enumerate}

\[
f (x) = \int_ {\mathrm{H}} f (x h) \varrho (h) y d \nu (y)
\]

where \(\varphi\in\mathcal{K}(\mathrm{G})\) and \(y\) belongs to the space of \(\varrho\) .)

\begin{enumerate}
\def\labelenumi{\arabic{enumi})}
\setcounter{enumi}{11}
\item
  Let G = R and let \(\pi\) be the regular representation of G in \(\mathrm{L}^{2}(\mathbf{R})\) . Prove that \(\chi \leqslant \pi\) for every character \(\chi \in \widehat{G}\) .
\item
  Let \(\varrho\) be an irreducible unitary representation of G in a Hilbert space E, and \(\pi\) a unitary representation of G in F. Let \(x \in E\) be of norm 1. One assumes that \(\varrho \leqslant \pi\) (cf.~Exercise 11).
\end{enumerate}

We denote by \(A \subset \mathcal{C}(G)\) the set of normalized diagonal matrix coefficients of \(\pi\) , and \(\widetilde{A}\) the closed convex hull of A in \(L^{\infty}(G)\) for the weak topology.

Let \(x \in E\) be of norm 1 and let \(f: g \mapsto \langle x \mid \varrho(g)x \rangle\) be the associated diagonal matrix coefficient.

\begin{enumerate}
\def\labelenumi{\alph{enumi})}
\item
  The set \(\widetilde{\mathrm{A}}\) is a compact convex set in \(\mathrm{L}^{\infty}(\mathrm{G})\) and \(f\) is an extreme point of \(\widetilde{\mathrm{A}}\) .
\item
  The function f belongs to the closure of A in \(\widetilde{A}\) for the weak topology.
\item
  The function f belongs to the closure of A in \(\mathcal{C}(G)\) endowed with the topology of compact convergence.
\item
  The trivial representation of G in C is weakly contained in \(\pi\) if and only if, for every compact subset C of G and every real \(\varepsilon > 0\) , there exist \(x \in F\) of norm 1 such that
\end{enumerate}

\[
\sup _ {g \in \mathrm{C}} \| \pi (g) x - x \| <   \varepsilon .
\]

\begin{enumerate}
\def\labelenumi{\arabic{enumi})}
\setcounter{enumi}{13}
\tightlist
\item
  Let G be a locally compact group. One says that G has the property ]T[ of Kazhdan \(^{(1)}\) if a unitary representation \(\pi\) of G has a nonzero trivial subrepresentation \(^{(2)}\) if and only if the trivial representation \(\mathbf{1}_{\text{G}}\) of G on C is weakly contained in \(\pi\) .
\end{enumerate}

\begin{enumerate}
\def\labelenumi{\alph{enumi})}
\item
  Let G be a locally compact group having property ]T[. Let H be the set of open subgroups of G that are generated by a compact subset of G. The group G is the union of the subgroups H ∈ H.
\item
  Let \(\pi\) be the representation of G that is the Hilbert sum of the quasi-regular representations of G on G/H for \(H \in H\) (cf. exercise 2 of V, p. 486). We have \(1_{G} \leqslant \pi\) .
\item
  The group G is generated by a compact subset of G. In particular, if G is discrete, then G is generated by a finite subset of G.
\end{enumerate}

\begin{enumerate}
\def\labelenumi{\arabic{enumi})}
\setcounter{enumi}{14}
\tightlist
\item
  Let G be a locally compact group, C a subset of G, and \(\varepsilon > 0\) . One says that \((\mathrm{C}, \varepsilon)\) is a Kazhdan pair for G if, for every unitary representation \(\pi\) of G in a Hilbert space E, the existence of a vector \(x \in E\)
\end{enumerate}

such that

\[
\sup _ {g \in \mathrm{C}} \| \pi (g) x - x \| <   \varepsilon \| x \|,
\]

  implies the existence of a nonzero trivial subrepresentation of \(\pi\) . We say that C is a Kazhdan set if there exists \(\varepsilon > 0\) such that \((C, \varepsilon)\) is a Kazhdan pair.

\begin{enumerate}
\def\labelenumi{\alph{enumi})}
\tightlist
\item
  Let \((\mathrm{C}, \varepsilon)\) be a Kazhdan pair of G. Let \(\pi\) be a unitary representation of G in a Hilbert space E. Let \(x \in E\) be nonzero and \(\delta > 0\) such that
\end{enumerate}

\[
\sup _ {g \in \mathrm{C}} \| \pi (g) x - x \| <   \delta \varepsilon \| x \|.
\]

We then have \(\|p(x)-x\|\leqslant\delta\|x\|\) , where p is the orthogonal projector of E whose image is the space of G-invariant vectors.

\begin{enumerate}
\def\labelenumi{\alph{enumi})}
\setcounter{enumi}{1}
\item
  A locally compact group G has property ]T[ if and only if there exists a Kazhdan pair (C,ε). (To prove that property ]T[ implies the assertion, proceed by contraposition by considering the unitary representation that is the Hilbert sum of representations π \(_{C,\varepsilon}\) parametrized by a compact subset C of G and ε > 0 such that π \(_{C,\varepsilon}\) contains no nontrivial subrepresentation but there exists x \(_{C,\varepsilon}\) of norm 1 with sup \(_{g\in\mathbf{C}}\|\pi(g)x_{\mathrm{C},\varepsilon}-x_{\mathrm{C},\varepsilon}\|<\varepsilon.\)
\item
  Let G be a locally compact group having property ]T[. Every compact subset C of G that generates G is a Kazhdan set.
\item
  Let G be a locally compact group having property ]T[ and let C be a compact subset of G with nonempty interior that is a Kazhdan set. Then C generates G. In particular, if G is discrete, then the Kazhdan sets are precisely the finite generating subsets of G. (The subgroup H of G generated by C is open; consider the quasi-regular representation \(\pi\) of G on G/H and show that the projection of the characteristic function \(\varphi\) of H \(\in\) G/H onto the space of invariant vectors of \(\pi\) is equal to \(\varphi\) .)
\end{enumerate}

\begin{enumerate}
\def\labelenumi{\arabic{enumi})}
\setcounter{enumi}{15}
\item
  Let G be a locally compact group that is countable at infinity and let H be a closed subgroup of G. If G has property ]T[, and if G/H admits a bounded G-invariant measure, then H has property ]T[. (Use Exercise 20 of V, p.~492.)
\item
  Let \(n \geqslant 2\) be an integer. We consider the real Lie group \(\mathbf{G} = \mathbf{R}^n \rtimes \mathbf{SL}(n, \mathbf{R})\) , and we denote by H the subgroup \(\mathbf{R}^n \times \{e\}\) of \(\mathbf{G}\) , identified with \(\mathbf{R}^n\) .
\end{enumerate}

Let \(\pi\) be a unitary representation of G in a complex Hilbert space E. Suppose that for every compact subset C of G and every \(\varepsilon > 0\) , there exists \(x \in E\) nonzero such that

\[
\sup _ {g \in \mathrm{C}} \| \pi (g) x - x \| <   \varepsilon \| x \|.
\]

\begin{enumerate}
\def\labelenumi{\alph{enumi})}
\item
  There exists a sequence \((x_{m})\) of vectors of norm 1 in E such that the functions \(\varphi_{m}(g)=\langle x_{m}|\pi(g)x_{m}\rangle\) converge to the constant function 1 in \(\mathcal{C}(\mathrm{G})\) endowed with the topology of compact convergence. (Consider a sequence \((\mathrm{C}_{m})_{m\in\mathbf{N}}\) of compact subsets whose union is equal to G, and take \(x_{m}\) satisfying the above condition for \((\mathrm{C}_{m},1/(m+1))\) .)
\item
  For every \(m \in N\) , there exists a positive measure \(\mu_{m}\) of total mass 1 on \(R^{n}\) whose Fourier transform is identified with the restriction of \(\varphi_{m}\) to H.
\item
  Suppose that \(\mu_{m}(\{0\}) = 0\) for all \(m \in N\) . We then identify \(\mu_{m}\) with a positive measure of total mass 1 on \(R^{n} - \{0\}\) , and we denote by \(\nu_{m}\) the image measure of \(\mu_{m}\) on projective space \(\mathbf{P}_{n-1}(\mathbf{R})\) . There exists a subsequence \((\nu_{n_{k}})_{k \in \mathbf{N}}\) converging weakly to a positive measure \(\nu\) of total mass 1 on \(\mathbf{P}_{n-1}(\mathbf{R})\) ; the measure \(\nu\) is \(\mathbf{SL}(n, \mathbf{R})\) -invariant. Deduce a contradiction.
\item
  Let \(m \in N\) be such that \(\mu_{m}(\{0\}) > 0\) . Let \(\psi\) be the restriction of \(\varphi_{m}\) to H. There exists c > 0 such that \(\psi = c + \widetilde{\psi}\) where \(\widetilde{\psi}\) is a function of positive type on \(R^{n}\) .
\item
  The representation \(\pi\) contains H-invariant vectors. (Use Exercise 2.) One says that the pair (G,H) has the relative ]T[ property.
\end{enumerate}

\begin{enumerate}
\def\labelenumi{\arabic{enumi})}
\setcounter{enumi}{17}
\tightlist
\item
  Let \(\mathrm{G} = \mathbf{SL}(2, \mathbf{R})\) . Let N and P denote the subgroups
\end{enumerate}

\[
\mathrm{N} = \Big \{\left( \begin{array}{c c} 1 & t \\ 0 & 1 \end{array} \right) | t \in \mathbf {R} \Big \}, \qquad \mathrm{P} = \Big \{\left( \begin{array}{c c} a & t \\ 0 & a ^ {- 1} \end{array} \right) | a \in \mathbf {R} _ {*}, t \in \mathbf {R} \Big \}
\]

of G. Let \(\pi\) be a unitary representation of G in a complex Hilbert space E. Suppose that there exists \(x \in \mathrm{E}\) such that \(\pi(n)x = x\) for all \(n \in \mathrm{N}\) .

\begin{enumerate}
\def\labelenumi{\alph{enumi})}
\tightlist
\item
  The homogeneous space \(\mathbf{SL}(2, \mathbf{R}) / \mathrm{N}\) is identified with \(\mathbf{R}^2 - \{0\}\) by the mapping \(g \mapsto ge_1\) where \(e_1\) is the first vector of the canonical basis of \(\mathbf{R}^2\) , and the homogeneous space \(\mathbf{SL}(2, \mathbf{R}) / \mathrm{P}\) is identified with the subspace \((\mathbf{R} - \{0\}) \times \{0\}\) of \(\mathbf{SL}(2, \mathbf{R}) / \mathrm{N}\) .
\item
  Every continuous function \(f\) on G such that \(f(ngn') = f(g)\) for all \((n, n') \in \mathrm{N}^2\) satisfies \(f(pgp') = f(g)\) for all \((p, p') \in \mathrm{P}^2\) . (Interpret \(f\) as a continuous function \(\varphi\) on \(\mathbf{R}^2 - \{0\}\) invariant under the natural action of N and verify that it is constant on \(\mathbf{SL}(2, \mathbf{R}) / \mathrm{P}\) .)
\end{enumerate}

\begin{enumerate}
\def\labelenumi{\alph{enumi})}
\setcounter{enumi}{2}
\tightlist
\item
  One has \(\pi(g)x = x\) for all \(g \in G\) .
\end{enumerate}

\begin{enumerate}
\def\labelenumi{\arabic{enumi})}
\setcounter{enumi}{18}
\tightlist
\item
  Let \(n \geqslant 2\) be an integer. We consider the real Lie group \(\mathrm{G} = \mathbf{SL}(n, \mathbf{R})\) . Let N and H denote the closed subgroups of G whose elements are of block form
\end{enumerate}

\[
\left( \begin{array}{c c} \operatorname{Id} _ {n - 1} & x \\ 0 & 1 \end{array} \right), \quad (x \in \mathbf {R} ^ {n - 1}), \qquad \left( \begin{array}{c c} g & 0 \\ 0 & 1 \end{array} \right), \quad (g \in \mathbf {S L} (n - 1, \mathbf {R}))
\]

respectively. The group N is identified with \(R^{n-1}\) and the group H is identified with \(\mathbf{SL}(n-1,\mathbf{R})\) .

\begin{enumerate}
\def\labelenumi{\alph{enumi})}
\item
  Let \(\pi\) be a unitary representation of G in a complex Hilbert space E. An element \(x\) of E is invariant under G if and only if it is invariant under N.
\item
  If \(n \geqslant 3\) , then the group G has property ]T[. (Combine the result of the previous question with the previous exercise.)
\item
  If \(n \geqslant 3\) , then the group \(\widetilde{G} = R^{n} \rtimes G\) has property ]T[. (Let \(\pi\) be a unitary representation such that \(1_{\widetilde{G}} \leqslant \pi\) ; there exists a vector \(x \neq 0\) invariant under the subgroup G of \(\widetilde{G}\) by the preceding question; show that the corresponding matrix coefficient is constant.)
\item
  If \(n \geqslant 3\) , the group \(\mathbf{SL}(n, \mathbf{Z})\) has property ]T[. (Use Exercise 7 of INT, VII, p.~116.)
\item
  Let \(n \geqslant 3\) . The group \(\mathbf{SL}(n, \mathbf{Z})\) is finitely generated. Let H be the family of normal subgroups of finite index of \(\mathbf{SL}(n, \mathbf{Z})\) ; it is infinite. For every finite generating subset S of \(\mathbf{SL}(n, \mathbf{Z})\) the family of Cayley graphs \((\mathcal{G}(\mathbf{SL}(n, \mathbf{Z})/\mathrm{H}, \mathrm{SH}/\mathrm{H}))_{\mathrm{H} \in \mathcal{H}}\) is a family of expander graphs (cf.~TA, 2, p.~221, exercise 7 and exercise 16 of IV, p.~323).
\end{enumerate}

\begin{enumerate}
\def\labelenumi{\arabic{enumi})}
\setcounter{enumi}{19}
\tightlist
\item
  Let G be a locally compact group and \(\mu\) a left Haar measure on G. We denote by \(\Delta\) the modulus of G. For any function f on G, we will denote by \(\widetilde{f} = \Delta f^{*}\) .
\end{enumerate}

A measure \(\nu\) on G is said to be of positive type if \(\langle \nu, \tilde{f} * f \rangle \geqslant 0\) for every function \(f \in \mathcal{K}(\mathrm{G})\) . A locally integrable function \(\varphi\) on G is said to be of positive type if \(\langle \varphi, f^* * f \rangle \geqslant 0\) for every function \(f \in \mathcal{K}(\mathrm{G})\) .

\begin{enumerate}
\def\labelenumi{\alph{enumi})}
\item
  Suppose that \(\nu\) is bounded. The measure \(\nu\) is of positive type if and only if \(\gamma(\nu) \geqslant 0\) .
\item
  Let \(\varphi\in\mathcal{K}(\mathrm{G})\) . The function \(\varphi\) is of positive type as a locally integrable function if and only if it is so according to Definition 4 of V, p.~442.
\item
  Let \(\nu\) be a measure of positive type on G. The formula
\end{enumerate}

\[
(f \mid g) _ {\nu} = \langle \nu , \widetilde {f} * g \rangle
\]

defines a positive Hermitian form on \(\mathcal{K}(\mathrm{G})\) . We denote by \(\mathrm{E}_{\nu}\) the separated completion Hilbert space of \(\mathcal{K}(\mathrm{G})\) equipped with this Hermitian form.

\begin{enumerate}
\def\labelenumi{\alph{enumi})}
\setcounter{enumi}{3}
\tightlist
\item
  For \(f \in \mathcal{K}(\mathrm{G})\) , the function \(\sigma(g)f \colon g \mapsto (\varepsilon_g * f)\Delta^{1/2}\) belongs to \(\mathcal{K}(\mathrm{G})\) and \(\| \sigma(g)f \|_{\nu} = \|f\|_{\nu}\) . The map \(g \mapsto \sigma(g)\) extends to a unitary representation of G in \(\mathrm{E}_{\nu}\) .
\end{enumerate}

\begin{enumerate}
\def\labelenumi{\alph{enumi})}
\setcounter{enumi}{4}
\tightlist
\item
  Let U be an open neighborhood of e in G, and let \(\nu\) be a positive-type measure on G such that there exists a real number \(c \geqslant 0\) so that
\end{enumerate}

\[
\langle \nu , f * \widetilde {f} \rangle \leqslant c \left(\int_ {\mathrm{G}} f d \mu\right) ^ {2}
\]

  if \(f \in \mathcal{K}(\mathrm{G})\) is positive and vanishes outside U. There exists a continuous function of positive type \(\varphi\) on G such that \(\nu = \varphi \Delta^{-1/2} \cdot \mu\) . (Let \(j\) be the canonical morphism from \(\mathcal{K}(\mathrm{G})\) into \(\mathrm{E}_{\nu}\) ; if \(\mathfrak{B}\) is a basis for the filter of neighborhoods of \(e\) whose elements are contained in U and \(f_{\mathrm{V}} \in \mathcal{K}_{+}(\mathrm{G})\) is a function with support in \(V \in B\) of integral 1, prove that \(j(f_{\mathrm{V}})\) converges weakly according to B to an element \(y_{0}\) of \(E_{\nu}\) such that

\[
(j (f) \mid y _ {0}) = \int_ {\mathrm{G}} f ^ {*} d \mu
\]

for \(f \in \mathcal{K}(\mathrm{G})\) ; setting \(\varphi(g) = (y_0 \mid \sigma(g)y_0)\) , verify that \(\varphi\) has the required properties.)

\begin{enumerate}
\def\labelenumi{\alph{enumi})}
\setcounter{enumi}{5}
\tightlist
\item
  Let \(\varphi_{1}\) be a continuous positive-definite function on G and \(\varphi_{2}\) a locally integrable function of positive type. If \(\varphi_{1}-\varphi_{2}\) is of positive type, then \(\varphi_{2}\) is equal locally almost everywhere to a continuous function of positive type. (Verify that one can apply the previous question.)
\end{enumerate}

\begin{enumerate}
\def\labelenumi{\arabic{enumi})}
\setcounter{enumi}{20}
\tightlist
\item
  We retain the notation of the preceding exercise.
\end{enumerate}

\begin{enumerate}
\def\labelenumi{\alph{enumi})}
\item
  Let \(f \in \mathcal{L}^{2}(G)\) that is of positive type as a locally integrable function. The partial operator on \(L^{2}(G)\) with domain \(\mathcal{K}(G)\) defined by \(\varphi \mapsto \gamma(\varphi)f\) is positive. We denote by \(\varrho(f)\) the Friedrichs extension of f (example 8 of IV, p.~295); it is a positive self-adjoint partial operator on \(L^{2}(G)\) .
\item
  Let \(f \in \mathcal{L}^2(\mathrm{G})\) . One says that \(f\) is sober if there exists a real number \(c \geqslant 0\) so that \(\| \boldsymbol{\gamma}(\varphi)f \| \leqslant c \| \varphi \|\) for all \(\varphi \in \mathcal{K}(\mathrm{G}) \subset \mathcal{L}^2(\mathrm{G})\) . We then denote by \(\varrho(f)\) the endomorphism of \(\mathrm{L}^2(\mathrm{G})\) which extends \(\varphi \mapsto \boldsymbol{\gamma}(\varphi)f\) . This definition coincides with the previous one if \(f\) is of positive type, and this condition then amounts to saying that \(\varrho(f)\) is positive.
\item
  Let \(f_{1}\) and \(f_{2}\) be sober functions in \(\mathcal{L}^{2}(G)\) . If \(\varrho(f_{1}) = \varrho(f_{2})\) , then \(f_{1}\) and \(f_{2}\) are equal in \(L^{2}(G)\) .
\end{enumerate}

\begin{enumerate}
\def\labelenumi{\alph{enumi})}
\setcounter{enumi}{3}
\tightlist
\item
  Let \(f \in \mathcal{L}^2(\mathrm{G})\) be of positive type. For all \(g \in \mathrm{G}\) , we have \(\varrho(f) \circ \boldsymbol{\gamma}(g) = \boldsymbol{\gamma}(g) \circ \varrho(f)\) . Moreover, for every \(h \in \operatorname{dom}(\varrho(f))\) , we have \(\varrho(f)h = h * f\) .
\end{enumerate}

\begin{enumerate}
\def\labelenumi{\alph{enumi})}
\setcounter{enumi}{4}
\tightlist
\item
  Let \(f_{1}\) and \(f_{2}\) be sober elements of \(\mathcal{L}^{2}(\mathrm{G})\) such that \(\varrho(f_{1})\) and \(\varrho(f_{2})\) commute. The function \(f_{1}*f_{2}\) is a continuous positive-definite function on G; it is sober, and we have
\end{enumerate}

\[
\varrho (f _ {1} * f _ {2}) = \varrho (f _ {1}) \varrho (f _ {2}), \quad \langle f _ {1} | f _ {2} \rangle \geqslant 0.
\]

\begin{enumerate}
\def\labelenumi{\alph{enumi})}
\setcounter{enumi}{5}
\tightlist
\item
  With the notations and hypotheses of the previous question, if \(f_{2} - f_{1}\) is of positive type, then
\end{enumerate}

\[
\left\| f _ {2} - f _ {1} \right\| ^ {2} \leqslant \left\| f _ {2} \right\| ^ {2} - \left\| f _ {1} \right\| ^ {2}.
\]

\begin{enumerate}
\def\labelenumi{\alph{enumi})}
\setcounter{enumi}{6}
\tightlist
\item
  Let \((f_n)_{n\in \mathbf{N}}\) be a sequence of sober positive-type elements of \(\mathcal{L}^2 (\mathrm{G})\) , pairwise commuting, and such that \(f_{n + 1} - f_n\) is of positive type for all \(n\in \mathbf{N}\) . If the sequence \((\| f\| _n)\) is bounded, then the sequence \((f_n)\) converges in \(\mathcal{L}^2 (\mathrm{G})\) .
\end{enumerate}

\begin{enumerate}
\def\labelenumi{\arabic{enumi})}
\setcounter{enumi}{21}
\tightlist
\item
  We retain the notation of the preceding exercise.
\end{enumerate}

Let \(f\) be a continuous function of positive type belonging to \(\mathcal{L}^2 (\mathrm{G})\) .

\begin{enumerate}
\def\labelenumi{\alph{enumi})}
\item
  There exists an increasing sequence \((p_{n})_{n\in\mathbf{N}}\) of polynomial functions on [0,1] which converges uniformly on [0,1] to the function \(t\mapsto\sqrt{t}\) .
\item
  Suppose that \(f\) is sober and that \(0 \leqslant f \leqslant 1\) in \(\mathcal{L}(\mathrm{L}^2(\mathrm{G}))\) . There exists a sequence \((h_n)\) in \(\mathcal{L}^2(\mathrm{G})\) satisfying the following conditions:
\end{enumerate}

\begin{enumerate}
\def\labelenumi{(\roman{enumi})}
\item
  Each \(h_n\) is sober and of positive type;
\item
  The function \(h_{n+1} - h_n\) is of positive type for all \(n \in \mathbf{N}\) ;
\item
  The endomorphisms \(\varrho(h_n)\) commute pairwise;
\item
  One has \(\varrho(h_n)^2 \leqslant \varrho(f)\) for all \(n\) .
\end{enumerate}

\begin{enumerate}
\def\labelenumi{\alph{enumi})}
\setcounter{enumi}{2}
\tightlist
\item
  The sequence \((h_n)_{n\in \mathbf{N}}\) converges in \(\mathcal{L}^2 (\mathrm{G})\) ; if \(h\) is a limit of \((h_n)\) , then \(h\) is of positive type and sober, and satisfies \(\varrho (h)^2 = \varrho (f)\) ; we have \(f = h*h\) .
\end{enumerate}

\begin{enumerate}
\def\labelenumi{\alph{enumi})}
\setcounter{enumi}{3}
\tightlist
\item
  We no longer assume that \(f\) is sober. For \(t \in \mathbf{R}\) , we denote by \(p_t\) the spectral projection of the partial self-adjoint operator \(\varrho(f)\) defined by \([- \infty, t]\) . We have \(p_t \circ \boldsymbol{\gamma}(g) = \boldsymbol{\gamma}(g) \circ p_t\) for all \(t \in \mathbf{R}\) and \(g \in G\) .
\end{enumerate}

\begin{enumerate}
\def\labelenumi{\alph{enumi})}
\setcounter{enumi}{4}
\tightlist
\item
  For \(t \in \mathbf{R}\) , set \(f_t = p_t(f)\) ; this is the class in \(\mathrm{L}^2(\mathrm{G})\) of a continuous positive-definite sober function in \(\mathcal{L}^2(\mathrm{G})\) (Check that \(\gamma(\varphi)f_t = \varrho(f)p_t(g)\) for \(g \in \mathcal{H}(\mathrm{G})\) , and deduce that \(\varrho(f_t) = \varrho(f) \circ p_t\) , then apply the results of Exercise 20.)
\end{enumerate}

\begin{enumerate}
\def\labelenumi{\alph{enumi})}
\setcounter{enumi}{5}
\tightlist
\item
  There exists \(h \in \mathcal{L}^2(\mathrm{G})\) such that \(f = h * h\) (Consider \(f_n = p_n(f)\) for \(n \in \mathbf{N}\) and apply part \(c\) ).)
\end{enumerate}

\begin{enumerate}
\def\labelenumi{\arabic{enumi})}
\setcounter{enumi}{22}
\tightlist
\item
  We denote \(G = R_{+}^{*}\) ; it is a locally compact commutative group and the measure \(\mu = dx/x\) is a Haar measure on G. The group G and the group iR are in duality via the mapping \((x, it) \mapsto x^{it}\) ; the dual measure of \(\mu\) is Lebesgue measure on iR (cf. Cor. 4 of II, p. 236). We denote by F the Fourier transform of G, and \(\overline{F}\) the Fourier cotransform of \(\widehat{G}\) .
\end{enumerate}

We denote by X the identity function on G.

\begin{enumerate}
\def\labelenumi{\alph{enumi})}
\item
  Let \(\nu\) be a positive measure on G. We denote by \(\mathrm{I}_{\nu}\) the set of \(\sigma \in \mathbf{R}\) such that the function \(\mathrm{X}^{\sigma}\) is \(\nu\) -integrable, and \(\mathrm{B}_{\nu}\) the set of \(s \in \mathbf{R}\) such that \(\mathcal{R}(s) \in \mathrm{I}_{\nu}\) . The set \(\mathrm{I}_{\nu}\) is an interval of \(\mathbf{R}\) and the function \(s \mapsto \nu(\mathrm{X}^s)\) is defined and holomorphic in the interior of \(\mathrm{B}_{\nu}\) .
\item
  Let \(f\) be a locally integrable function on G. We denote by \(\mathrm{I}_{f} = \mathrm{I}_{f\cdot \mu}\) and \(\mathrm{B}_f = \mathrm{B}_{f\cdot \mu}\) . For every \(s \in \mathrm{B}_f\) we set
\end{enumerate}

\[
\widehat {f} (s) = \int_ {\mathrm{G}} f (x) x ^ {s} d \mu (x).
\]

If there exists \(\delta > 0\) and \(c \geqslant 0\) such that \(|\widehat{f}(s)| \leqslant c(1 + |s|)^{-1 - \delta}\) for \(s \in \mathrm{B}_f\) , then for all \(\sigma\) in the interior of \(\mathrm{B}_f\) , we have

\[
f (x) = \frac {1}{2 \pi} \int_ {\mathbf {R}} \widehat {f} (\sigma + i t) x ^ {- \sigma - i t} d t.
\]

\begin{enumerate}
\def\labelenumi{\alph{enumi})}
\setcounter{enumi}{2}
\tightlist
\item
  Let \(f \in \mathcal{L}^{2}(\mathrm{G})\) . We assume that 0 belongs to the interior of \(I_{f}\) . We then have \(\overline{\mathcal{F}} f(it) = \widehat{f}(it)\) for almost every \(t \in \mathbf{R}\) .
\end{enumerate}

\begin{enumerate}
\def\labelenumi{\alph{enumi})}
\setcounter{enumi}{3}
\item
  Let \(f\) and \(g\) be measurable functions on \(\mathbf{G}\) . Let \(\sigma\) be a real number such that \(\mathrm{X}^{-2\sigma}f\in \mathrm{L}^2 (\mathrm{G})\) and \(\mathrm{X}^{2\sigma}g\in \mathrm{L}^2 (\mathrm{G})\) . We then have

\[
\int_ {\mathrm{G}} f g d \mu = \frac {1}{2 \pi} \int_ {\mathbf {R}} \widehat {f} (- \sigma - i t) \widehat {g} (\sigma + i t) d t.
\]

\end{enumerate}

\begin{enumerate}
\def\labelenumi{\alph{enumi})}
\setcounter{enumi}{4}
\tightlist
\item
  Let \(f\) and \(g\) be measurable functions on \(G\) . One assumes that
\end{enumerate}

\begin{enumerate}
\def\labelenumi{(\roman{enumi})}
\item
  There exists \(c \in \mathbf{R}\) such that \(|f(x)| \leqslant cx / (1 + x)\) for all \(x \in \mathbf{G}\) ;
\item
  One has ]0,1[ ⊂ I \(_{g}\) and for every σ such that 0 < σ < 1, the function X \(^{\sigma}\) g is bounded;
\item
  There exists \(c' \in \mathbf{R}\) such that \(|\widehat{g}(s)| \leqslant c'(1 + |s|)^{-2}\) for all \(s \in \mathbf{C}\) such that \(\mathcal{R}(s) \in ]0, 1[\) .
\end{enumerate}

Then we have

\[
\int_ {\mathrm{G}} f g d \mu = \frac {1}{2 \pi} \int_ {\mathbf {R}} \widehat {f} (- \sigma - i t) \widehat {g} (\sigma + i t) d t
\]

for all \(\sigma\in\mathbf{R}\) such that \(0<\sigma<1\) .

\begin{enumerate}
\def\labelenumi{\alph{enumi})}
\setcounter{enumi}{5}
\tightlist
\item
  Let \(\sigma \in \mathbf{R}\) and let \(f \in \mathcal{C}(\sigma + i\mathbf{R})\) . There exists a positive measure \(\nu\) on G such that \(\sigma \in \mathrm{I}_{\nu}\) and \(\widehat{\nu}(\sigma + it) = f(\sigma + it)\) for all \(t \in \mathbf{R}\) if and only if for every integer \(n \geqslant 1\) , every family \((t_i)_{0 \leqslant i \leqslant n}\) of complex numbers and any family \((z_i)_{0 \leqslant i \leqslant n}\) of complex numbers such that \(\mathcal{R}(z_i) = \sigma\) for all \(i\) , we have
\end{enumerate}

\[
\sum_ {i = 0} ^ {n} \sum_ {j = 0} ^ {n} \bar {t} _ {i} t _ {j} f \left(\frac {1}{2} (\overline {{z}} _ {i} + z _ {j})\right) \geqslant 0.
\]

